% answers_short.tex —— 2010 年·全国新课标卷·理科 · 简短答案速查
% 由 tools/make_short_answers.py 自 answers.tex 的 \gksolhead/\gkanswer 提取，
% 并与 2010/2010_数学真题_全国新课标理_答案要点.md 逐题比对后生成，勿手改（改源 MD 后重跑）。

\begin{gknotice}
\textbf{使用说明}：本节只给最终答案便于快速核对，完整推导见解析版。答案为编写者独立推导，\textbf{非官方标准答案}。
\end{gknotice}

\gkpart{选择题}
\begin{center}\begin{tabular}{ccccccccccccc}
\toprule
题号 & 1 & 2 & 3 & 4 & 5 & 6 & 7 & 8 & 9 & 10 & 11 & 12 \\
\midrule
答案 & D & C & A & C & C & B & D & B & A & B & C & B \\
\bottomrule
\end{tabular}\end{center}

\gkpart{填空题}
\begin{itemize}[leftmargin=2.6em]
  \item[\textbf{13.}] $\dfrac{N_1}{N}$
  \item[\textbf{14.}] ①②⑤
  \item[\textbf{15.}] $(x-3)^2+y^2=2$
  \item[\textbf{16.}] $60^\circ$
\end{itemize}

\gkpart{解答题}
\begin{itemize}[leftmargin=2.6em]
  \item[\textbf{17.}] (1) $a_n=2^{2n-1}$；(2) $S_n=\dfrac{2\left[1+(3n-1)4^n\right]}9$
  \item[\textbf{18.}] (1) 见解析版；(2) $\dfrac{\sqrt2}{4}$
  \item[\textbf{19.}] (1) $0.14$；(2) 有 $99\%$ 的把握认为是否需要帮助与性别有关；(3) 按性别分层随机抽样，再按总体比例加权估计
  \item[\textbf{20.}] (1) $e=\dfrac{\sqrt2}{2}$；(2) $\dfrac{x^2}{18}+\dfrac{y^2}{9}=1$
  \item[\textbf{21.}] (1) 减区间为 $(-\infty,0)$，增区间为 $(0,+\infty)$；(2) $a\leq\dfrac12$
  \item[\textbf{22.}] (1) 见解析版；(2) 见解析版
  \item[\textbf{23.}] (1) $(1,0)$、$\left(\dfrac12,-\dfrac{\sqrt3}{2}\right)$；(2) $x=\dfrac12\sin^2\alpha$，$y=-\dfrac12\sin\alpha\cos\alpha$；轨迹为 $\left(x-\dfrac14\right)^2+y^2=\dfrac1{16}$
  \item[\textbf{24.}] (1) 图象为以 $(2,1)$ 为顶点、左右斜率分别为 $-2\text{、}2$ 的折线；(2) $a\in\left[\dfrac12,+\infty\right)$
\end{itemize}
